% ============ 3 平展基本群 (原书 pp.25-30) ============
\chapter{平展基本群}\label{sec:3}
\index{平展基本群}
\index{有限平展覆盖}
\index{Galois 覆盖}
\index{通用覆盖空间}
\index{Riemann 存在定理}
\index{几何点}

平展基本群分类簇(或概形)的有限平展覆盖，方式与通常的基本群分类拓扑空间的覆盖空间相同.本节只是该主题的综述---完整叙述见本节末尾的参考文献.

我们先从函子观点回顾经典理论.

\subsection*{拓扑基本群}

设 $X$ 是连通拓扑空间.为了有一个好的理论，我们假定 $X$ 道路连通且半局部单连通(即每个 $P \in X$ 都有邻域 $U$，使得 $U$ 中基于 $P$ 的每个闭路都可在 $X$ 中收缩到 $P$).

固定 $x \in X$.基本群 $\pi_1(X, x)$ 定义为 $X$ 中基于 $x$ 的闭路的同伦类的群.

连续映射 $\varphi \colon Y \to X$ 是 $X$ 的覆盖空间，如果 $X$ 的每个点 $P$ 都有开邻域 $U$，使得 $\varphi^{-1}(U)$ 是开集 $U_i$ 的不交并，其中每个 $U_i$ 都被 $\varphi$ 同胚地映到 $U$ 上.覆盖空间的映射 $(Y, \varphi) \to (Y', \varphi')$ 是满足 $\varphi' \circ \alpha = \varphi$ 的连续映射 $\alpha \colon Y \to Y'$.在我们的假设下，存在单连通覆盖空间 $\widetilde{\varphi} \colon \widetilde{X} \to X$.固定映到 $x \in X$ 的 $\tilde{x} \in \widetilde{X}$.则 $(\widetilde{X}, \tilde{x})$ 有如下泛性质：对任意覆盖空间 $Y \to X$ 与映到 $x \in X$ 的 $y \in Y$，存在唯一的覆盖空间映射 $\widetilde{X} \to Y$ 把 $\tilde{x}$ 送到 $y$.特别地，对 $(\widetilde{X}, \tilde{x})$，它唯一到唯一的同构；它称为 $(X, x)$ 的通用覆盖空间.

设 $\operatorname{Aut}_X(\widetilde{X})$ 表示覆盖空间映射 $\widetilde{X} \to \widetilde{X}$ 的群，$\alpha \in \operatorname{Aut}_X(\widetilde{X})$.因为 $\alpha$ 是覆盖空间映射，$\alpha \tilde{x}$ 也映到 $x \in X$.因此，从 $\tilde{x}$ 到 $\alpha \tilde{x}$ 的道路被 $\widetilde{\varphi}$ 映为 $X$ 中基于 $x$ 的闭路.因为 $\widetilde{X}$ 单连通，该闭路的同伦类不依赖于道路的选择，于是以这种方式得到映射 $\operatorname{Aut}_X(\widetilde{X}) \to \pi_1(X, x)$.这个映射是同构.

上述论断的证明见 Greenberg, M., \emph{Lectures on Algebraic Topology}, Benjamin, 1966 第 I 部分.

设 $\mathrm{Cov}(X)$ 是 $X$ 的只有有限多个连通分支的覆盖空间的范畴---态射是覆盖空间映射.我们定义 $F \colon \mathrm{Cov}(X) \to \mathbf{Sets}$ 为把覆盖空间 $\varphi \colon Y \to X$ 送到集合 $\varphi^{-1}(x)$ 的函子.这个函子被 $\widetilde{X}$ 表示，即
\[
F(Y) \cong \operatorname{Hom}_X(\widetilde{X}, Y) \quad \text{关于 } Y \text{ 函子性地成立}.
\]
确实，我们上面注意到，给出覆盖空间映射 $\widetilde{X} \to Y$ 与给出点 $y \in \varphi^{-1}(x)$ 是一回事.

若让 $\operatorname{Aut}_X(\widetilde{X})$ 右作用于 $\widetilde{X}$，则它左作用于 $\operatorname{Hom}_X(\widetilde{X}, Y)$：
\[
\alpha f \overset{\mathrm{def}}{=} f \circ \alpha; \qquad \alpha \in \operatorname{Aut}_X(\widetilde{X}), \ f \colon \widetilde{X} \to Y.
\]
于是，我们看到 $F$ 可以看作从 $\mathrm{Cov}(X)$ 到 $\operatorname{Aut}_X(\widetilde{X})$(或 $\pi_1(X, x)$)集合范畴的函子.

$\pi_1(X, x)$ 分类 $X$ 的覆盖空间这一事实被如下论断优美地概括：函子 $F$ 定义了从 $\mathrm{Cov}(X)$ 到只有有限多个轨道的 $\pi_1(X, x)$ 集合范畴的等价.

\subsection*{平展基本群}

设 $X$ 是连通簇(或概形).我们选取几何点 $\bar{x} \to X$，即坐标在可分代数闭域中的 $X$ 的点.当 $X$ 是代数闭域 $k$ 上的簇时，可取 $\bar{x}$ 为 $X(k)$ 的元素.对概形 $X$，选取 $\bar{x}$ 等同于选取点 $x \in X$ 连同包含剩余域 $\kappa(x)$ 的可分代数闭域.

我们不再能把基本群定义为闭路的同伦类构成的群.例如，若 $X$ 是 $\mathbb{C}$ 上的光滑曲线，则 $X(\mathbb{C})$ 是 Riemann 曲面，$X(\mathbb{C})$ 上的闭路有复(或代数)维数 $\frac{1}{2}$---我们不是物理学家，不允许这样的对象.作为替代，我们模仿把基本群刻画为通用覆盖空间的覆盖变换群.

回忆有限平展映射 $\varphi \colon Y \to X$ 是开的(因为平展---见命题~\ref{prop:2.12})且闭的(因为有限---见 AG 8.7)，从而它是满射(只要 $Y \neq \varnothing$).若 $X$ 是代数闭域上的簇，$\varphi \colon Y \to X$ 有限且平展，则 $\varphi$ 的每个纤维都有完全相同数目的点.此外，每个 $x \in X$ 都有平展邻域 $(U, u) \to (X, x)$，使得 $Y \times_{X} U$ 是开子簇(或子概形) $U_i$ 的不交并，其中每个 $U_i$ 被 $\varphi \times 1$ 同构地映到 $U$ 上(见后面).因此，有限平展映射是有限覆盖空间的自然类比.

我们定义 $\mathbf{FEt}/X$ 为如下范畴：对象是有限平展映射 $\varphi \colon Y \to X$(有时称为 $X$ 的有限平展覆盖)，箭头是 $X$-态射.

定义 $F \colon \mathbf{FEt}/X \to \mathbf{Sets}$ 为把 $(Y, \varphi)$ 送到 $Y$ 的位于 $x$ 之上的 $\bar{x}$ 值点的集合的函子，即 $F(Y) = \operatorname{Hom}_X(\bar{x}, Y)$.若 $X$ 是代数闭域上的簇且 $\bar{x} \in X(k)$，则 $F(Y) = \varphi^{-1}(\bar{x})$.

我们希望把 $X$ 的通用覆盖空间定义为表示 $F$ 的对象，但不幸的是，(通常)不存在这样的对象.例如，设 $\mathbb{A}^{1}$ 是特征零的代数闭域 $k$ 上的仿射线.则 $\mathbb{A}^{1} \setminus \{0\}$ 的有限平展覆盖是映射
\[
t \mapsto t^{n} \colon \mathbb{A}^{1} \setminus \{0\} \to \mathbb{A}^{1} \setminus \{0\}
\]
(见例~\ref{ex:2.6}).在这些覆盖中，没有``最大的''一个---在拓扑情形($k = \mathbb{C}$)，通用覆盖是
\[
\mathbb{C} \xrightarrow{\exp} \mathbb{C} \setminus \{0\},
\]
它没有代数对应物.

然而，函子 $F$ 是 pro-可表示的.这意味着存在 $X$ 的有限平展覆盖的射影系统 $\widetilde{X} = (X_i)_{i \in I}$，它以有向集\footnote{有向集是集合 $A$ 连同自反，传递的二元关系 $\leq$，使得对所有 $a, b \in A$，存在 $c \in A$ 满足 $a \leq c$ 且 $b \leq c$.}$I$ 为指标集，使得
\[
F(Y) = \operatorname{Hom}(\widetilde{X}, Y) \overset{\mathrm{def}}{=} \varinjlim_{i \in I} \operatorname{Hom}(X_i, Y) \quad \text{关于 } Y \text{ 函子性地成立}.
\]
在上一段考虑的例子中，$\widetilde{X}$ 是以按整除偏序的正整数为指标集的族 $t \mapsto t^{n} \colon \mathbb{A}^{1} \setminus \{0\} \to \mathbb{A}^{1} \setminus \{0\}$.

我们称 $\widetilde{X}$ 为 $X$ 的``那个''通用覆盖空间.可以选取 $\widetilde{X}$ 使每个 $X_i$ 在 $X$ 上都是 Galois 的，即它在 $X$ 上的次数等于 $\operatorname{Aut}_X(X_i)$ 的阶.映射 $X_j \to X_i$($i \leq j$)诱导同态 $\operatorname{Aut}_X(X_j) \to \operatorname{Aut}_X(X_i)$，我们定义
\[
\pi_1(X, \bar{x}) = \operatorname{Aut}_X(\widetilde{X}) \overset{\mathrm{def}}{=} \varprojlim_{i} \operatorname{Aut}_X(X_i),
\]
赋予其作为有限离散群的射影极限的自然拓扑.若 $X_n \to \mathbb{A}^{1} \setminus \{0\}$ 表示上一段中的覆盖，则
\[
\operatorname{Aut}_X(X_n) = \mu_n(k) \quad (\text{$k$ 中 $1$ 的 $n$ 次单位根群}),
\]
其中 $\zeta \in \mu_n(k)$ 通过 $x \mapsto \zeta x$ 作用.于是
\[
\pi_1(\mathbb{A}^{1} \setminus 0, \bar{x}) = \varprojlim_{n} \mu_n(k) \cong \widehat{\mathbb{Z}}.
\]
这里 $\widehat{\mathbb{Z}} \cong \prod_{\ell} \mathbb{Z}_{\ell}$ 是 $\mathbb{Z}$ 对由有限指数子群定义的拓扑的完备化.该同构通过选取同构 $Z/n\mathbb{Z} \to \mu_n(k)$ 的相容系来定义，或者等价地，对每个 $n$ 选取 $1$ 的本原 $n$ 次单位根 $\zeta_n$，使得对所有 $m, n > 0$ 有 $\zeta^{m}_{mn} = \zeta_n$.

$\pi_1(X, \bar{x})$ 在 $\widetilde{X}$ 上的(右)作用为 $X$ 的每个有限平展覆盖 $Y$ 定义了 $\pi_1(X, \bar{x})$ 在 $F(Y)$ 上的左作用.当 $F(Y)$ 带离散拓扑时，这个作用是连续的---这无非是说它穿过某个 $i \in I$ 的有限商 $\operatorname{Aut}_X(X_i)$.

\begin{theorem}\label{thm:3.1}
函子 $Y \mapsto F(Y)$ 定义了从 $X$ 的有限平展覆盖范畴到有限离散 $\pi_1(X, \bar{x})$ 集合范畴的等价.
\end{theorem}

于是，$\pi_1(X, \bar{x})$ 分类 $X$ 的有限平展覆盖，方式与拓扑基本群分类拓扑空间的覆盖空间相同.

\begin{example}\label{ex:3.2}
设 $\mathbb{P}^{1}$ 表示代数闭域上的射影直线.$\mathbb{P}^{1}$ 上的微分 $\omega = dz$ 除 $\infty$ 处外没有极点，而在 $\infty$ 处有二重极点.对任意连通且有限平展的覆盖 $\varphi \colon Y \to \mathbb{P}^{1}$，$\varphi^{*}\omega$ 在位于 $\infty$ 之上的 $\deg(\varphi)$ 个点的每一个处都有二重极点，且没有其他极点.因此 $\varphi^{*}\omega$ 的因子有次数 $-2\deg(\varphi)$.由于这等于 $2\operatorname{genus}(Y) - 2$，$\deg(\varphi)$ 必须等于 1，从而 $\varphi$ 是同构.因此 $\pi_1(\mathbb{P}^{1}, \bar{x}) = 1$，正如预期.

同样的论证表明，当基域特征为零时，$\pi_1(\mathbb{A}^{1}, \bar{x}) = 0$.更一般地，不存在连通曲线 $Y$ 与次数 $> 1$ 的有限映射 $Y \to \mathbb{P}^{1}$，使其在 $\mathbb{A}^{1}$ 上非分歧且在 $\infty$ 上驯分歧.

$\mathbb{P}^{n}$ 的平展基本群也为零.
\end{example}

\begin{remark}\label{rem:3.3}
(a) 若 $\bar{\bar{x}}$ 是 $X$ 的第二个几何点，则存在同构 $\pi_1(X, \bar{x}) \to \pi_1(X, \bar{\bar{x}})$，它良定义到共轭.

(b) 基本群 $\pi_1(X, \bar{x})$ 是 $(X, \bar{x})$ 的共变函子.

(c) 设 $k \subset \Omega$ 是特征零的代数闭域.对 $k$ 上的任意簇 $X$，函子 $Y \mapsto Y_{\Omega} \colon \mathbf{FEt}/X \to \mathbf{FEt}/X_{\Omega}$ 是范畴等价，从而对 $X_{\Omega}$ 的任意几何点 $\bar{x}$，定义同构 $\pi_1(X, \bar{x}) \cong \pi_1(X_{\Omega}, \bar{x})$.对特征 $p$ 中的 $p$-覆盖，这个论断是假的；例如，对每个 $\alpha \in k[T]$，我们有 $\mathbb{A}^{1}$ 的由方程
\[
Y^{p} - Y = \alpha
\]
定义的 Artin-Schreier 有限平展覆盖，当把 $k$ 换成更大的域时得到更多的这种覆盖.
\end{remark}

\subsection*{\texorpdfstring{$\mathbb{C}$ 上的簇}{C 上的簇}}

设 $X$ 是 $\mathbb{C}$ 上的簇，$Y \to X$ 是 $X$ 的有限平展覆盖.若 $X$ 非奇异，则 $Y$ 也非奇异，并且当 $X(\mathbb{C})$ 与 $Y(\mathbb{C})$ 赋予复拓扑时，$Y(\mathbb{C}) \to X(\mathbb{C})$ 成为 $X(\mathbb{C})$ 的有限覆盖空间.

让我们能把经典基本群与平展基本群联系起来的关键结果是：

\begin{theorem}\label{thm:3.4}[Riemann 存在定理]
设 $X$ 是 $\mathbb{C}$ 上的非奇异簇.把 $X$ 的有限平展覆盖 $(Y, \varphi)$ 送到 $X(\mathbb{C})$ 的有限覆盖空间 $(Y(\mathbb{C}), \varphi)$ 的函子是范畴等价.
\end{theorem}

证明的困难部分是证明该函子本质满，即 $X(\mathbb{C})$ 的每个有限(拓扑)覆盖空间都有自然的代数结构.证明见：Grauert 和 Remmert, Math. Ann. (1958), 245--318; SGA 4, XI.4.3；以及 SGA 1, XII.下面第 21 节包含最后那个证明的梗概.

由该定理推出，$X$ 的平展通用覆盖空间 $\widetilde{X} = (X_i)_{i \in I}$ 有性质：$X(\mathbb{C})$ 的每个有限拓扑覆盖空间都是某个 $X_i(\mathbb{C})$ 的商.设 $x = \bar{x}$ 是 $X(\mathbb{C})$ 的任意元素.则
\[
\pi_1(X, x) \overset{\mathrm{def}}{=} \varprojlim_{i} \operatorname{Aut}_X(X_i) = \varprojlim_{i} \operatorname{Aut}_{X(\mathbb{C})}(X_i(\mathbb{C})) = \widehat{\pi_1(X(\mathbb{C}), x)},
\]
其中 $\widehat{\pi_1(X(\mathbb{C}), x)}$ 是 $\pi_1(X(\mathbb{C}), x)$ 对由有限指数子群定义的拓扑的完备化.

由于 $\pi_1(\mathbb{C} \setminus \{0\}, x) \cong \mathbb{Z}$，我们恢复事实 $\pi_1(\mathbb{A}^{1} \setminus \{0\}, \bar{x}) \cong \widehat{\mathbb{Z}}$.

\begin{remark}\label{rem:3.5}
设 $X$ 是特征零的代数闭域 $k$ 上的非奇异簇，$\sigma$ 与 $\tau$ 是 $k$ 到 $\mathbb{C}$ 的嵌入.则对平展基本群，
\[
\pi_1(X, \bar{x}) \cong \pi_1({}^{\sigma}X, {}^{\sigma}\bar{x}) \cong \pi_1({}^{\tau}X, {}^{\tau}\bar{x})
\]
(由注记~\ref{rem:3.3}(c)).因此
\[
\widehat{\pi_1(({}^{\sigma}X)(\mathbb{C}), {}^{\sigma}\bar{x})} \cong \widehat{\pi_1(({}^{\tau}X)(\mathbb{C}), {}^{\tau}\bar{x})}.
\]
正如我们在引言中注意到的，存在这样的例子：
\[
\pi_1(({}^{\sigma}X)(\mathbb{C}), {}^{\sigma}\bar{x}) \not\cong \pi_1(({}^{\tau}X)(\mathbb{C}), {}^{\tau}\bar{x}).
\]
本节中的论断(适当修改后)对 $\mathbb{C}$ 上有限型的每个连通概形 $X$ 成立；特别地，对这样的概形，$\pi_1(X, \bar{x}) \cong \widehat{\pi_1(X(\mathbb{C}), x)}$.
\end{remark}

\subsection*{例}

\textbf{域的谱} 对 $X = \operatorname{Spec} k$($k$ 是域)，平展态射 $Y \to X$ 是平展 $k$-代数的谱，且每个都是有限的.因此，与其用 $\mathbf{FEt}/X$，不如用平展 $k$-代数的反范畴 $\mathbf{Et}/k$.

$X$ 的几何点的选取等同于包含 $k$ 的可分代数闭域 $\Omega$ 的选取.定义 $F \colon \mathbf{Et}/k \to \mathbf{Sets}$ 为
\[
F(A) = \operatorname{Hom}_k(A, \Omega).
\]
设 $\tilde{k} = (k_i)_{i \in I}$ 是包含在 $\Omega$ 中的所有有限 Galois 扩张组成的射影系统.则 $\tilde{k}$ ind-表示 $F$，即
\[
F(A) \cong \operatorname{Hom}_k(A, \tilde{k}) \overset{\mathrm{def}}{=} \varinjlim_{i \in I} \operatorname{Hom}_k(A, k_i) \quad \text{关于 } A \text{ 函子性地成立},
\]
---这是显然的.定义
\[
\operatorname{Aut}_k(\tilde{k}) = \varprojlim_{i \in I} \operatorname{Aut}_{k\text{-alg}}(k_i).
\]
于是
\[
\operatorname{Aut}_k(\tilde{k}) = \varprojlim_{i} \operatorname{Gal}(k_i/k) = \operatorname{Gal}(k^{\mathrm{sep}}/k),
\]
其中 $k^{\mathrm{sep}}$ 是 $k$ 在 $\Omega$ 中的可分代数闭包.此外，$F$ 定义了从 $\mathbf{Et}/k$ 到有限离散 $\operatorname{Gal}(k^{\mathrm{sep}}/k)$ 集合范畴的范畴等价.这个论断概括了通常的域的 Galois 理论，并且容易由它推出.

\textbf{正规簇(或概形)} 对连通正规簇(或概形) $X$，最自然的是取位于 $X$ 的一般点 $x$ 之上的几何点 $\bar{x}$.当然，严格地说，若 $X$ 是簇，我们不能这么做，因为簇没有一般点，但这等同于选取包含 $X$ 上有理函数域 $k(X)$ 的可分代数闭域 $\Omega$.设 $L$ 是 $\Omega$ 中 $k(X)$ 的所有有限可分域扩张 $K$ 的并，其中 $X$ 在 $K$ 中的正规化在 $X$ 上平展；则
\[
\pi_1(X, \bar{x}) = \operatorname{Gal}(L/k(X))
\]
带 Krull 拓扑.

\textbf{Albanese 簇} [待补充.]

\textbf{计算平展基本群} 对 $\mathbb{C}$ 上的连通簇，我们可以利用与经典基本群的关系来计算平展基本群.对特征零的任意代数闭域上的连通簇，我们可以利用如下事实：从一个代数闭域到另一个代数闭域作基变换时基本群不变.

对特征 $p \neq 0$ 的代数闭域 $k$ 上的连通簇 $X_0$，有时可以通过把 $X_0$ 提升到特征零来获得 $\pi_1(X_0, x)$ 的``非 $p$''部分的信息.假定存在完备离散赋值环 $R$，其剩余域为 $k$，分式域 $K$ 特征为零，以及 $R$ 上光滑且紧合的概形 $X$，其特殊纤维为 $X_0$.存在典范满同态 $\pi_1(X_{K^{\mathrm{al}}}, \bar{x}) \to \pi_1(X_0, \bar{x}_0)$，其核包含在 $\pi_1(X_{K^{\mathrm{al}}}, \bar{x})$ 到阶与 $p$ 互素的有限群的每个同态的核中；特别地，这两个群有相同的阶与 $p$ 互素的有限商.由于每条光滑射影曲线 $X_0$ 都能提升到特征零，得到 $\pi_1(X_0, \bar{x}_0)$(忽略 $p$-商)的描述：它是带标准关系的 $2\operatorname{genus}(X_0)$ 个生成元的自由群的完备化.\footnote{Grothendieck 对特征 $p$ 中次数与 $p$ 互素的曲线的有限平展覆盖的分类是他最早的重大成功之一，说明了他的代数几何方法的重要性.据我所知，这个结果至今没有纯代数的证明.}

设 $X$ 是域 $k$ 上的连通代数簇，并假定 $X_{k^{\mathrm{sep}}}$ 也连通.则存在正合序列
\[
1 \to \pi_1(X_{k^{\mathrm{sep}}}, \bar{x}) \to \pi_1(X, \bar{x}) \to \operatorname{Gal}(k^{\mathrm{sep}}/k) \to 1.
\]
这些映射来自基本群的函子性.

\textbf{数学中最有趣的对象} 可以说，它是
\[
\pi_1(X, \bar{x}),
\]
其中 $X$ 是 $\mathbb{Q}$(而不是 $\mathbb{C}$!)上挖去三个点 $0, 1, \infty$ 的射影直线.上述正合序列在这种情况下变为：
\[
1 \to \pi_1(X_{\mathbb{Q}^{\mathrm{al}}}, \bar{x}) \to \pi_1(X, x) \to \operatorname{Gal}(\mathbb{Q}^{\mathrm{al}}/\mathbb{Q}) \to 1.
\]
这里 $\pi_1(X_{\mathbb{Q}^{\mathrm{al}}}, \bar{x})$ 典范同构于挖去三个点的 Riemann 球面的基本群的完备化，该基本群由围绕三个被挖去点的闭路 $\gamma_0, \gamma_1, \gamma_\infty$ 生成，带有唯一关系 $\gamma_0\gamma_1\gamma_\infty = 1$.群 $\operatorname{Gal}(\mathbb{Q}^{\mathrm{al}}/\mathbb{Q})$ 已经很神秘，而理解这个扩张涉及混合 Tate 动机，数域的 $K$ 群，$\zeta$ 函数的特殊值等.见 Deligne 的文章，\emph{Galois Groups over Q}, Springer, 1989, pp 79--297.

显然 Grothendieck 是第一个强调这个群的意义的人---见 \emph{Geometric Galois Actions}, 1. \emph{Around Grothendieck's Esquisse d'un Programme},(Schneps 和 Lochak 编), Cambridge University Press, 1997.还有 Y. Ihara,G. Anderson,Minhyong Kim 等人的若干篇论文.

\begin{aside}\label{asd:3.6}
对正规代数簇，容易定义 $\pi_1$：它是函数域的最大非分歧扩张的 Galois 群.然而，Grothendieck 的方法允许你变动基点，定义基本群胚，等等.
\end{aside}

\subsection*{参考文献}

原始文献是 Grothendieck 的讨论班 SGA 1.最有用的叙述是：
\begin{quote}
  Murre, J.P., \emph{Lectures on an Introduction to Grothendieck's Theory of the Fundamental Group}, Tata Institute, 1967.
\end{quote}
另见：
\begin{quote}
  Szamuely, Tam\'as. \emph{Galois groups and fundamental groups}. Cambridge Studies in Advanced Mathematics, 117. Cambridge University Press, Cambridge, 2009.
\end{quote}
